\documentclass[./main.tex]{subfiles}
\begin{document}

% Slide # 10 
\begin{frame}[t, label=slide10]
        % Title
        \frametitle{Functional framework}

        % Body
        \footnotesize
        
        An observed sample path can be thought of as a sequence $\{X_{n}\}_{n=1}^{N}\subset I \subseteq \mathbb{R}$ of $N>0$ i.i.d. random variables sampled from $p(x)\in H:=\{f\in L^{2}(\mathbb{R})\;:\;\int_{\mathbb{R}}^{}f(x)dx = 1\}$.

        \vspace{0.25cm}
        \uncover<2->{
                Fixing $h>0$ identifies a unique partition $I_{h}$ of $I$ into $0 < M < N$ homogeneous (i.e. fixed width), disjoint subintervals
                \begin{equation*}
                        I_{h} = \cup_{k=1}^{M}B_{k}\,,\quad B_{k} = \big[\min\{X_{t}\} + (k-1)h, \min\{X_{t}\} + kh\big)\,.
                \end{equation*}
        }

        \uncover<3->{
                Construct the subspace of $H$ of piecewise constant functions
                \begin{equation*}
                        \mathcal{H}_{h}:=\bigg\{f\in H\;:\;f(x)=\sum_{k=1}^{M}c_{k}\boldsymbol{1}_{B_{k}}(x)\bigg\}\,.
                \end{equation*}
        }

        \uncover<4->{
                \begin{proposition}
                     There exists a unique minimiser of
                     \begin{equation*}
                         p_{h}:=\text{argmin}_{f\in \mathcal{H}_{h}}||f - p||_{H}\,,
                     \end{equation*}
                     and it coincides with the bin-averaged true density, i.e. $c_{k}=\frac{1}{h}\int_{B_{k}}^{}p(x)dx=:p_{k}$.
                \end{proposition}
        }

\end{frame}

\end{document}
